% GENERATED FILE. Edit canonical lessons, then run npm run generate:books.
% canonical-source-hash: fnv1a64:e93ff0c293ca4bf5
% canonical-lessons: FA-C202-pare, FA-C202-kaj, FA-C202-khoshtip, FA-C202-kamrang, FA-C202-bimaze

\chapter{Torn, Crooked, Stylish, Pale, Bland}
\label{ch:fa-torn-crooked-stylish-pale-bland}
\hlchaptermodality{202}

\begin{chapteropening}
\textbf{By the end of this chapter:} \emph{I can say torn, crooked, stylish, pale and bland.}

Complete the last lesson of chapter 202: I can say torn, crooked, stylish, pale and bland.
\end{chapteropening}

\section[\texorpdfstring{\fa{پاره} (pâre)}{پاره (pâre)}]{\fa{پاره} (pâre) — torn}

\label{lesson:FA-C202-pare}



\begin{quote}
Before the new one: say the Persian for cloudy, then the Persian for rainy.
\end{quote}

\subsection*{You'll want to know: \fa{پاره}}
\textbf{\fa{پاره}} — \emph{pâre} — \textquotedblleft{}torn\textquotedblright{}.

\begin{cousinweb}
\emph{pâre} is a native Persian word, inherited rather than borrowed; it is usually traced to Middle Persian \emph{pārag}, \textquotedblleft{}piece\textquotedblright{}. Persian uses it for \textquotedblleft{}torn, in pieces\textquotedblright{}.
\end{cousinweb}

\begin{grammarlens}[title={What you've built}]
One more word for this chapter.
\end{grammarlens}

\subsection*{Your turn}
\begin{itemize}
\raggedright
  \item \emph{Say it:} \emph{pâre}
  \item \emph{Say it:} \emph{pâre}, once more
  \item \emph{Say it:} say \emph{pâre}
  \item \emph{Recall:} say the Persian for cloudy, then the Persian for rainy, then say \emph{pâre} again
\end{itemize}

\begin{tcolorbox}[breakable,colback=teal!4,colframe=teal!35!black,title={Before you move on}]
What does \fa{پاره} mean? (\textquotedblleft{}Torn\textquotedblright{}.) Say it once more.
\end{tcolorbox}
\section[\texorpdfstring{\fa{کج} (kaj)}{کج (kaj)}]{\fa{کج} (kaj) — crooked}

\label{lesson:FA-C202-kaj}



\begin{quote}
Before the new one: say the Persian for rainy, then the Persian for torn.
\end{quote}

\subsection*{You'll want to know: \fa{کج}}
\textbf{\fa{کج}} — \emph{kaj} — \textquotedblleft{}crooked\textquotedblright{}.

\begin{cousinweb}
\emph{kaj} is a native Persian word, inherited rather than borrowed; it is usually traced to Middle Persian \emph{kaž}, \textquotedblleft{}crooked\textquotedblright{}.
\end{cousinweb}

\begin{grammarlens}[title={What you've built}]
One more word for this chapter.
\end{grammarlens}

\subsection*{Your turn}
\begin{itemize}
\raggedright
  \item \emph{Say it:} \emph{kaj}
  \item \emph{Say it:} \emph{kaj}, once more
  \item \emph{Say it:} say \emph{kaj}
  \item \emph{Recall:} say the Persian for rainy, then the Persian for torn, then say \emph{kaj} again
\end{itemize}

\begin{tcolorbox}[breakable,colback=teal!4,colframe=teal!35!black,title={Before you move on}]
What does \fa{کج} mean? (\textquotedblleft{}Crooked\textquotedblright{}.) Say it once more.
\end{tcolorbox}
\section[\texorpdfstring{\fa{خوشتیپ} (khoshtip)}{خوشتیپ (khoshtip)}]{\fa{خوشتیپ} (khoshtip) — stylish, good-looking}

\label{lesson:FA-C202-khoshtip}



\begin{quote}
Before the new one: say the Persian for torn, then the Persian for crooked.
\end{quote}

\subsection*{You'll want to know: \fa{خوشتیپ}}
\textbf{\fa{خوشتیپ}} — \emph{khoshtip} — \textquotedblleft{}stylish, good-looking\textquotedblright{}.

\begin{cousinweb}
\emph{khoshtip} is \emph{khosh}, \textquotedblleft{}pleasant, good\textquotedblright{}, with \emph{tip}, \textquotedblleft{}type, look\textquotedblright{}, borrowed from French \emph{type}. A Persian-French hybrid.
\end{cousinweb}

\begin{grammarlens}[title={What you've built}]
One more word for this chapter.
\end{grammarlens}

\subsection*{Your turn}
\begin{itemize}
\raggedright
  \item \emph{Say it:} \emph{khoshtip}
  \item \emph{Say it:} \emph{khoshtip}, once more
  \item \emph{Say it:} say \emph{khoshtip}
  \item \emph{Recall:} say the Persian for torn, then the Persian for crooked, then say \emph{khoshtip} again
\end{itemize}

\begin{tcolorbox}[breakable,colback=teal!4,colframe=teal!35!black,title={Before you move on}]
What does \fa{خوشتیپ} mean? (\textquotedblleft{}Stylish, good-looking\textquotedblright{}.) Say it once more.
\end{tcolorbox}
\section[\texorpdfstring{\fa{کمرنگ} (kamrang)}{کمرنگ (kamrang)}]{\fa{کمرنگ} (kamrang) — pale (of colour); weak (of tea)}

\label{lesson:FA-C202-kamrang}



\begin{quote}
Before the new one: say the Persian for crooked, then the Persian for stylish.
\end{quote}

\subsection*{You'll want to know: \fa{کمرنگ}}
\textbf{\fa{کمرنگ}} — \emph{kamrang} — \textquotedblleft{}pale (of colour); weak (of tea)\textquotedblright{}.

\begin{cousinweb}
\emph{kamrang} is \emph{kam}, \textquotedblleft{}little\textquotedblright{}, with \emph{rang}, \textquotedblleft{}colour\textquotedblright{}: little colour. Both are native; \emph{rang} continues Middle Persian \emph{rang} (EtymDB).
\end{cousinweb}

\begin{grammarlens}[title={What you've built}]
One more word for this chapter.
\end{grammarlens}

\subsection*{Your turn}
\begin{itemize}
\raggedright
  \item \emph{Say it:} \emph{kamrang}
  \item \emph{Say it:} \emph{kamrang}, once more
  \item \emph{Say it:} say \emph{kamrang}
  \item \emph{Recall:} say the Persian for crooked, then the Persian for stylish, then say \emph{kamrang} again
\end{itemize}

\begin{tcolorbox}[breakable,colback=teal!4,colframe=teal!35!black,title={Before you move on}]
What does \fa{کمرنگ} mean? (\textquotedblleft{}Pale (of colour); weak (of tea)\textquotedblright{}.) Say it once more.
\end{tcolorbox}
\section[\texorpdfstring{\fa{بیمزه} (bimaze)}{بیمزه (bimaze)}]{\fa{بیمزه} (bimaze) — bland, tasteless}

\label{lesson:FA-C202-bimaze}



\begin{quote}
Before the new one: say the Persian for stylish, then the Persian for pale.
\end{quote}

\subsection*{You'll want to know: \fa{بیمزه}}
\textbf{\fa{بیمزه}} — \emph{bimaze} — \textquotedblleft{}bland, tasteless\textquotedblright{}.

\begin{cousinweb}
\emph{bimaze} is built from \emph{bi-}, \textquotedblleft{}without\textquotedblright{}, and \emph{maze}, \textquotedblleft{}taste\textquotedblright{}: tasteless. Both are native Persian words, inherited rather than borrowed.
\end{cousinweb}

\begin{grammarlens}[title={What you've built}]
One more word for this chapter.
\end{grammarlens}

\subsection*{Your turn}
\begin{itemize}
\raggedright
  \item \emph{Say it:} \emph{bimaze}
  \item \emph{Say it:} \emph{bimaze}, once more
  \item \emph{Say it:} say \emph{bimaze}
  \item \emph{Recall:} say the Persian for stylish, then the Persian for pale, then say \emph{bimaze} again
\end{itemize}

\begin{tcolorbox}[breakable,colback=teal!4,colframe=teal!35!black,title={Before you move on}]
What does \fa{بیمزه} mean? (\textquotedblleft{}Bland, tasteless\textquotedblright{}.) Say it once more.
\end{tcolorbox}
